% ============================================================
%  Chapter 7 — Discussion
% ============================================================
\section{Discussion}
\label{ch:discussion}

\subsection{Regulatory Implications}
\label{sec:disc:regulatory}

The results of the six scenarios reveal a nuanced picture of how the
Carbon Intensity Indicator regulation interacts with the operational reality
of a fixed-schedule Baltic Sea service.  Three overarching regulatory
implications emerge.

\paragraph{Schedule dominates CII compliance in the near term.}
Under the 2026 Rating~C threshold, the published timetable --- not the
\CII{} constraint --- determines achievable speeds and emissions.  The
33.9\,\% compliance slack observed in S3 implies that the vessel already
operates well within the current regulatory ceiling, and that speed
optimisation alone cannot further reduce emissions without renegotiating
port windows.  This finding challenges the assumption embedded in much of
the IMO's slow-steaming policy narrative: that tightening the \CII{}
target will straightforwardly induce operators to reduce speed.  For
fixed-schedule short-sea services, the instrument is partially blunt.

\paragraph{The 2030 Rating~B target is structurally unachievable.}
S4 demonstrates that a 31\,\% reduction in required \CII{} relative to
the 2008 reference --- corresponding to a Rating~B target of
9.9291~\gGTnm{} --- cannot be met on the current timetable regardless of
speed choices.  The minimum achievable \CII{} under the schedule is
10.3187~\gGTnm{}, a gap of 0.39~\gGTnm{} (3.9\,\%).  This finding has
direct policy relevance: operators in analogous fixed-schedule trades cannot
rely on speed management alone to comply with aggressive post-2026 targets.
The regulatory pathway from Rating~C (2026) to Rating~B (2030) requires
complementary interventions such as schedule redesign, alternative fuels,
or energy efficiency measures.

\paragraph{Carbon pricing materially increases voyage cost but does not
change the optimal speed profile.}
In S4b, the EU \ETS{} charge (EUR~100/t\,CO$_2$) adds a significant
surcharge to the voyage cost without altering the optimal speed solution,
because the schedule --- not the emission cost --- remains binding.  This
implies that, for fixed-schedule operators, carbon pricing exerts a
transfer-payment effect (shifting cost from operators to carbon permit
markets) rather than a behavioural effect (inducing emission reductions
through speed change).  The emission-reduction channel of EU \ETS{} is
therefore partially ineffective for this vessel type and service pattern
under the current timetable.  This is consistent with findings in the
broader literature on carbon pricing for captive transport routes
\parencite{kontovas2011}.

\subsection{Value of Schedule Flexibility}
\label{sec:disc:flexibility}

S5 quantifies the operational and economic benefit of replacing rigid
per-arc time windows with a global cycle-time cap.  The reduction in
attained \CII{} from 10.2705~\gGTnm{} (S3) to 8.7116~\gGTnm{} (S5)
--- a 15.2\,\% improvement --- reflects the ability to operate bottleneck
arcs at lower speeds when not constrained by individual port windows.

From a managerial perspective, this suggests that \emph{negotiations
with port authorities and cargo customers over arrival window flexibility
may yield greater emission reductions than any on-board technological
investment}.  This insight is particularly actionable for \RoRo{}
operators whose schedules are customer-driven: modest relaxation of
port-window tolerances could enable a step-change in \CII{} performance
without capital expenditure.

The value of flexibility also has implications for the design of future
\CII{} regulation.  If the IMO or EU were to provide explicit compliance
incentives for operators who negotiate flexible schedules (analogous to
the existing \textit{correction factor} mechanism for ice-class
operations), this could unlock emission reductions in the short-sea
sector that are currently inaccessible through speed management alone.

\subsection{Fuel Price Sensitivity}
\label{sec:disc:fuelprice}

HFO prices are highly volatile: in the period 2020--2025 they ranged from
approximately EUR~300/t to over EUR~900/t, driven by crude oil markets,
refinery margins, and the IMO~2020 sulphur cap.  This volatility raises
the question of how sensitive the operational conclusions are to the assumed
fuel price.

\paragraph{Speed profiles are invariant to fuel price in fixed-schedule scenarios.}
For scenarios with per-arc time budgets (S3, S4b), the minimum feasible
speed on each arc is determined solely by the distance--time relationship:
\begin{equation}
  v_{ij}^{\min} = \frac{d_{ij}}{B_{ij}}
  \label{eq:min_speed}
\end{equation}
Since the objective function minimises cost (fuel plus carbon) and the
fuel consumption is strictly increasing in speed, the cost-minimising
choice is always the minimum feasible speed on each arc, \emph{regardless
of the fuel price level}.  Changing the fuel price from EUR~400/t to
EUR~1{,}000/t does not alter any speed decision in S3 or S4b; it only
rescales the total voyage cost linearly.

This result has a strong policy implication.  The \CII{} regulation is designed
to induce operators to reduce speed by tightening the carbon intensity ceiling.
But for fixed-schedule operators, the schedule already forces minimum speeds
that are sufficient for \CII{} compliance --- and those minimum speeds are
optimal at any positive fuel price.  The regulatory instrument therefore
provides no marginal behavioural incentive.  Similarly, fuel price escalation
provides no additional compliance benefit: the vessel cannot slow further
than the schedule permits, and it would slow to that minimum anyway on
economic grounds.

\paragraph{Fuel price sensitivity of total cost.}
Although speeds are invariant, total voyage cost is sensitive to fuel price.
\Cref{tab:fuel_sensitivity} reports per-loop costs for S3, S4b, and S5
at four representative fuel prices, holding the \ETS{} carbon prices at
their scenario-specific values (EUR~40/t for 2026 scenarios; EUR~100/t
for S4b).

\begin{table}[h]
\centering
\caption{Fuel price sensitivity --- per-loop total cost (EUR)}
\label{tab:fuel_sensitivity}
\begin{tabular}{lrrrr}
  \toprule
  Scenario & \multicolumn{4}{c}{Fuel price (EUR/t HFO)} \\
  \cmidrule(lr){2-5}
  & 400 & 600 & 800 & 1{,}000 \\
  \midrule
  S3  (2026, fixed schedule)   & 465{,}639 & 643{,}174 & 820{,}709 & 998{,}245 \\
  S4b (2030, fixed schedule)   & 631{,}492 & 809{,}028 & 986{,}563 & 1{,}164{,}098 \\
  S5  (2026, flexible schedule)& 394{,}963 & 545{,}551 & 696{,}140 & 846{,}728 \\
  \midrule
  \textit{S5 saving vs.\ S3}   & \textit{70{,}676} & \textit{97{,}623} & \textit{124{,}569} & \textit{151{,}517} \\
  \bottomrule
\end{tabular}
\smallskip\\
{\footnotesize \ETS{} carbon cost held fixed: S3/S5 at EUR~40/t\,CO$_2$;
S4b at EUR~100/t\,CO$_2$.  Speed profiles are identical across all fuel
price levels within each scenario.
S4b values adjust only for the fuel price component relative to the
EUR~914/t baseline (CO$_2$ cost is EUR~276{,}422 for all S4b rows).}
\end{table}

\paragraph{The value of schedule flexibility increases with fuel price.}
The savings from moving to a flexible schedule (S5 versus S3) grow
proportionally with the fuel price, because S5 uses 134.7~t less fuel per
loop.  At EUR~400/t the annual saving is approximately EUR~1.20~million;
at EUR~1{,}000/t it rises to EUR~2.58~million --- a 2.1$\times$ amplification
over the same price range (EUR~70{,}676 to EUR~151{,}517 per loop).  This means that the business case for negotiating
relaxed port windows with cargo customers is strongest precisely when fuel
markets are expensive and volatile --- the conditions that characterise the
current market environment.  Operators facing sustained high oil prices should
therefore prioritise schedule flexibility discussions as the highest-return
\CII{} and cost management lever available without capital investment.

\subsection{Sensitivity to the Fuel Model}
\label{sec:disc:fuelmodel}

The thesis relies on the SCA polynomial fuel model, which exhibits a
local non-monotonicity at 12--13~kn in the light-load condition.  Two
implications follow.

First, the optimizer may select 12~kn over 13~kn in some scenarios
because the fuel rate at 12~kn is marginally lower.  In practice, this
difference is within the uncertainty of the polynomial fit and should not
be interpreted as a robust speed recommendation.  A sensitivity analysis
varying the polynomial coefficients would be prudent in future work.

Second, the non-monotonicity means that the fuel consumption curve does
not satisfy the usual convexity assumption exploited in continuous
relaxations of speed-optimisation models.  The present formulation
avoids this issue by discretising the speed set to integer knots and
solving a \MILP{} directly, but practitioners using derivative-based
methods should be aware of this feature of the SCA vessel's fuel map.

\subsection{Cargo and Loading Condition}
\label{sec:disc:cargo}

In the \RoRo{} \CII{} formula, cargo mass does not appear: the denominator
uses gross tonnage (\GT{}) rather than a transport work metric.  This
distinguishes the \RoRo{} case sharply from bulk and container shipping,
where utilisation directly affects the attained \CII{}.

Despite this, cargo loading condition affects the \emph{fuel rate}
(via the polynomial coefficients), and the per-arc cargo weights
observed in the EU MRV data vary substantially (light, medium, full).
The model assigns the mean cargo weight per arc to each scenario; this
may understate or overstate fuel consumption on individual legs.  A
stochastic extension treating cargo as a random variable would better
capture the distributional effects, though it would substantially
increase model complexity.

It is also worth noting that the effective capacity factor $f_m = 1.05$
applied in the CII denominator reflects the full-load capacity of the
vessel.  Under partial loading, the numerator (CO$_2$) decreases while
the denominator remains fixed, so the CII improves even though the
vessel is not efficiently loaded.  This is a feature of the regulatory
formula that may not accurately reflect the social cost of emissions per
unit of transport work delivered.

\subsection{Assumptions and Limitations}
\label{sec:disc:limitations}

Several modelling assumptions limit the generalisability of the results.

\paragraph{Post-2026 reduction factors.}
The 31\,\% reduction factor applied to derive the 2030 required \CII{}
is a scenario assumption, not a confirmed regulatory figure.  As of the
thesis writing date, the \IMO{} MEPC has adopted an 11\,\% reduction for
2026 but has not finalised post-2026 factors.  The 31\,\% assumption is
informed by the IMO 2023 GHG Strategy's 40\,\% fleet-wide ambition
\parencite{mepc377}, but the actual factor applicable to this vessel class
may differ.  If the eventual factor is lower (say, 20\,\%), the S4
infeasibility gap narrows or disappears; if higher (say, 40\,\%), the
gap widens and even Rating~C compliance requires schedule interventions.

\paragraph{Single vessel and route.}
The case study covers a single vessel (M/V Obbola) on a single Baltic
Sea route operated by \SCA{} Logistics.  The findings may not generalise
to other \RoRo{} vessels, different route structures, or operators with
greater schedule flexibility.  The bottleneck arc phenomenon is a
product of the specific port sequence and SCA's published timetable; a
vessel with shorter inter-port distances or more generous time windows
would not exhibit the same infeasibility.

\paragraph{Static voyage cost model.}
The model treats fuel prices and carbon prices as deterministic
parameters.  In reality, both are volatile: HFO prices have ranged from
EUR~400/t to over EUR~700/t in recent years, and EU \ETS{} allowance
prices have fluctuated between EUR~25/t and EUR~100/t.  A stochastic
programming extension that models price uncertainty would provide
more robust operational recommendations.

\paragraph{No technological alternatives.}
The model optimises speed within the existing vessel and fuel
configuration.  It does not consider alternative fuels (LNG, methanol,
ammonia), wind-assist technologies, or energy efficiency modifications
(hull coatings, propeller optimisation).  These interventions could
shift the entire \CII{} curve downward, potentially closing the S4
infeasibility gap.  Modelling the cost-effectiveness of these options
under the \CII{} regulatory pathway is a natural extension.

\paragraph{No berthing or port-congestion dynamics.}
The schedule model uses fixed time budgets derived from mean EU MRV
observations.  It does not model variable port waiting times, berth
capacity, tidal windows, or weather routing.  Incorporating these
would require a stochastic schedule model that is beyond the scope
of this thesis but represents an important direction for future work.

\subsection{Managerial Insights}
\label{sec:disc:managerial}

For operators of fixed-schedule \RoRo{} services on the Baltic Sea, the
following actionable insights emerge from this study.

\textbf{Near-term (2026):} \CII{} compliance is not a pressing concern
under the current regulatory trajectory and timetable.  The primary
operational lever is schedule adherence, not speed reduction.  No
significant investment in speed management is warranted until post-2026
targets are finalised.

\textbf{Medium-term (2030):} If post-2026 reduction factors approach
the IMO's 40\,\% fleet-wide ambition, even a Rating~C target may require
schedule renegotiation on bottleneck legs.  Operators should initiate
dialogue with port authorities and cargo customers now about the
feasibility of modest arrival-window relaxations.

\textbf{Carbon cost management:} The EU \ETS{} represents a growing
cost exposure that is not offset by speed reduction on fixed schedules.
Operators should budget for carbon costs as a fuel-proportional surcharge
and explore hedging strategies.

\textbf{Data infrastructure:} The validation exercise (S1 vs.\ EU MRV)
highlights the value of high-quality per-leg fuel consumption data.
The 1.2\,\% \CII{} error in S1 is acceptable, but improving the accuracy
of the polynomial fuel model --- particularly for the non-monotonic
region at 12--13~kn --- would benefit future scenario analyses.
